\chapter{Evaluation and Discussion}

In this chapter, we will take a closer look at the results of our work. We will examine the findings, limitations, and shortcomings of our study, as well as ways to address them. Finally, we will touch on future research in this area.

\section{Results}


During the course of the research, the file\_vault component was adapted for use by Snapper as a snapshot storage system. The integration was seamless for Snapper and other components; only minor changes were required to ensure proper functionality, changes that would have been necessary even without our component. Only a few modifications to lwext4 can be considered truly necessary specifically for this project. A comparative analysis of the resulting and target components can be seen in Table~\ref{tab:vault_comparison}.

\begin{table}[H]
\centering
\caption{Comparison of file\_vault and se\_vault}
\label{tab:vault_comparison}
\begin{tabular}{|p{4.2cm}|p{5.2cm}|p{5.2cm}|}
\toprule
\textbf{Parameter} & \textbf{file\_vault} & \textbf{se\_vault} \\
\midrule
Purpose & General-purpose secure file storage & Snapshot encryption for Phantom OS \\
\midrule
User interface & GUI via ui\_config / ui\_report ROM & None (headless) \\
\midrule
Authentication & Passphrase entered at runtime & Key file read from USB drive \\
\midrule
File system driver & rump\_vfs (ext2) & lwext4 (ext4) \\
\midrule
fsck support & No & Yes (configurable via \texttt{fsck\_apply}) \\
\midrule
Storage backend & Image file via File\_system session & Direct block device session \\
\midrule
Trust Anchor storage & File\_system session & USB block session \\
\midrule
Online extend / rekey & Yes & Removed \\
\midrule
Image file creation & Pre-allocated via truncate\_file & Skipped \\
\midrule
Block size & 512 & Standardized to 4~KB \\
\midrule
Init time (4~GB FS) & $\approx$20~min & $\approx$1~min \\
\bottomrule
\end{tabular}
\end{table}

Although we managed to secure the snapshots, performance has plummeted. This may not be as critical for less demanding tasks, since snapshots do not block Phantom processes for extended periods of time. However, during active operation, a large number of pages will be modified while a new snapshot is being created, and these will need to be written to the new snapshot, which in turn will lead to an even greater increase in snapshot creation time. This will also result in longer intervals between snapper calls, creating a risk of data loss over a longer period in the event of a failure. On the other hand, Dmitry Zavalishin described a problem where a large amount of memory is used and becomes obsolete between snapshot creations, citing graph computation as an example. When using the OS for similar tasks where a relatively long interval between snapshots is required and data loss during that period would not be critical, it is possible to mitigate the consequences to a significant extent. This is best verified through practical use, but that is currently impossible due to the system’s lack of readiness.

\section{Limitations}

This project certainly has a number of serious limitations. First, Tresor is not intended for production use. The author himself acknowledges this. The problem is that it is the only library designed for block-level encryption. And after the author left the team, its support has deteriorated even further. Second, testing was conducted on small amounts of user space. It cannot be said with certainty that the system will operate with comparable levels of performance and stability as the volume increases exponentially. Third, testing was conducted only on a virtual machine with placeholders instead of many actual working components. Although this scenario provides an indication of approximate results, it still differs significantly from real-world usage.

\section{Shortcomings}

The performance issue is obvious, as has been pointed out many times. The only apparent solution is to make corrections to the library. Another approach is to sacrifice security to some extent—though the exact extent remains to be determined—and implement a simpler encryption method. This can be done at the file system level. Genode includes an OpenSSL port, so there is no need to implement cryptographic algorithms from scratch, which would be unacceptable.

\section{Future work}

Clearly, we should focus on improving performance in the future. However, there is another aspect that hasn’t been discussed yet. It’s worth noting that isomem also uses a block device. This is part of the persistent storage where pages are swapped out from RAM. When the system shuts down, a large number of artifacts remain in it, allowing the machine’s state at the time of shutdown to be almost completely restored. This should be avoided. It is not difficult to modify se\_vault to provide a block session instead of a filesystem session. This will allow it to be used for encrypting the contents of the isomem block device. Moving the file system outside the sandbox will increase the solution’s flexibility. For greater security, conversely, one can place snapper and/or isomem inside the sandbox. This is a fairly simple solution to implement. In the early stages of adaptation, this option was also used in our work, but we abandoned it in favor of configurability. If it is absolutely necessary to protect data, it is still possible to do so using the \textit{block} VFS plugin, but given the current performance levels, this is undesirable. Another approach could be to attempt to combine snapper and tresor. Tresor also stores snapshots and uses incremental updates when writing. In this regard, their functionality is similar, but the components use it differently.
